\documentclass[10pt,a4paper]{amsart}
\usepackage{amsmath,amssymb,amsthm,graphicx}
\usepackage{verbatim}
\newenvironment{reference}{\comment}{\endcomment}
%\newenvironment{reference}{}{}
\newenvironment{slogan}{\comment}{\endcomment}
\newenvironment{history}{\comment}{\endcomment}

% For commutative diagrams we use Xy-pic
\usepackage[all]{xy}

% We use 2cell for 2-commutative diagrams.
\xyoption{2cell}
% Theorem environments.
%
\theoremstyle{plain}
\newtheorem{theorem}[subsection]{Theorem}
\newtheorem{proposition}[subsection]{Proposition}
\newtheorem{lemma}[subsection]{Lemma}

\theoremstyle{definition}
\newtheorem{definition}[subsection]{Definition}
\newtheorem{example}[subsection]{Example}
\newtheorem{exercise}[subsection]{Exercise}
\newtheorem{situation}[subsection]{Situation}

\theoremstyle{remark}
\newtheorem{remark}[subsection]{Remark}
\newtheorem{remarks}[subsection]{Remarks}

\numberwithin{equation}{subsection}

% Macros
%
\def\lim{\mathop{\mathrm{lim}}\nolimits}
\def\colim{\mathop{\mathrm{colim}}\nolimits}
\def\Spec{\mathop{\mathrm{Spec}}}
\def\Hom{\mathop{\mathrm{Hom}}\nolimits}
\def\Ext{\mathop{\mathrm{Ext}}\nolimits}
\def\SheafHom{\mathop{\mathcal{H}\!\mathit{om}}\nolimits}
\def\SheafExt{\mathop{\mathcal{E}\!\mathit{xt}}\nolimits}
\def\Sch{\mathit{Sch}}
\def\Mor{\mathop{\mathrm{Mor}}\nolimits}
\def\Ob{\mathop{\mathrm{Ob}}\nolimits}
\def\Sh{\mathop{\mathit{Sh}}\nolimits}
\def\NL{\mathop{N\!L}\nolimits}
\def\CH{\mathop{\mathrm{CH}}\nolimits}
\def\proetale{{pro\text{-}\acute{e}tale}}
\def\etale{{\acute{e}tale}}
\def\QCoh{\mathit{QCoh}}
\def\Ker{\mathop{\mathrm{Ker}}}
\def\Im{\mathop{\mathrm{Im}}}
\def\Coker{\mathop{\mathrm{Coker}}}
\def\Coim{\mathop{\mathrm{Coim}}}

% Boxtimes
%
\DeclareMathSymbol{\boxtimes}{\mathbin}{AMSa}{"02}

%
% Macros for moduli stacks/spaces
%
\def\QCohstack{\mathcal{QC}\!\mathit{oh}}
\def\Cohstack{\mathcal{C}\!\mathit{oh}}
\def\Spacesstack{\mathcal{S}\!\mathit{paces}}
\def\Quotfunctor{\mathrm{Quot}}
\def\Hilbfunctor{\mathrm{Hilb}}
\def\Curvesstack{\mathcal{C}\!\mathit{urves}}
\def\Polarizedstack{\mathcal{P}\!\mathit{olarized}}
\def\Complexesstack{\mathcal{C}\!\mathit{omplexes}}
% \Pic is the operator that assigns to X its picard group, usage \Pic(X)
% \Picardstack_{X/B} denotes the Picard stack of X over B
% \Picardfunctor_{X/B} denotes the Picard functor of X over B
\def\Pic{\mathop{\mathrm{Pic}}\nolimits}
\def\Picardstack{\mathcal{P}\!\mathit{ic}}
\def\Picardfunctor{\mathrm{Pic}}
\def\Deformationcategory{\mathcal{D}\!\mathit{ef}}
\title{The associated groupoid in sets inherits conditions (S1) and (S2)}
\author{The Stacks Project Authors}
\date{stacks-project:formal-defos.tex:local-snapshot-20261012; GNU FDL 1.2 or later}
\begin{document}
\maketitle
\noindent\textit{Editorial note.} All delivered mathematical text is the author's own, copied from the pinned Stacks Project chapter file formal-defos.tex; no statement or proof step was rewritten, reordered or removed.\par\medskip
\noindent\textit{Editorial note.} The counted proof is the author's complete argument: the transfer of (S1) through isomorphism classes, the construction of the two commutative diagrams over the square of rings encoding the second half of (S2), and the removal of the centralizer discrepancy using the three listed hypotheses.\par\medskip
\noindent\textit{Editorial note.} The label lemma-S1-S2-associated-functor names the counted statement, and the proof environment immediately following it in the source is delivered verbatim as its proof.\par\medskip
\noindent\textit{Editorial note.} The counted statement is carried by the author's own lemma environment in the chapter source.\par\medskip
\noindent\textit{Editorial note.} Pre-open gate: python operations/agent-authoring/tools/precard\_gate.py --source formal-defos.tex --label lemma-S1-S2-associated-functor --start 3094 --end 3188 --local-id 28014 returned RESULT=OK (L28014 free, occupancy FREE, public interval clear). Candidate taken from sources/stacks-formal-defos/candidates-halfA.json.\par\medskip
\noindent\textit{Editorial note.} The statement and proof are copied from the pinned chapter file; the cited inputs are kept as external pointers with their canonical Stacks tags.\par\medskip
\noindent\textit{Allowed standard prerequisites (editorial).}\begin{itemize}
\item Construction of the associated category fibred in sets. Same chapter, the construction taking a category cofibred in groupoids to the associated category fibred in sets by applying the set of isomorphism classes. Supplies the category F-bar whose objects and morphisms are the isomorphism classes of those of F.
\item Axioms of a category cofibred in groupoids. Standard consequences of the axioms of a category cofibred in groupoids, allowing the choice of the displayed commutative diagrams lying over a given square of rings. Supplies the diagrams exhibiting y over A x\_k k[epsilon] together with its pushforwards.
\end{itemize}\medskip
\begin{lemma}
\label{lemma-S1-S2-associated-functor}
Let $\mathcal{F}$ be a category cofibered in groupoids over
$\mathcal{C}_\Lambda$.
\begin{enumerate}
\item If $\mathcal{F}$ satisfies (S1), then so does
$\overline{\mathcal{F}}$.
\item If $\mathcal{F}$ satisfies (S2), then so does
$\overline{\mathcal{F}}$ provided at least one of the following conditions is
satisfied
\begin{enumerate}
\item $\mathcal{F}$ is a predeformation category,
\item the category $\mathcal{F}(k)$ is a set or a setoid, or
\item for any morphism $x_\epsilon \to x_0$ of $\mathcal{F}$
lying over $k[\epsilon] \to k$ the pushforward map
$\text{Aut}_{k[\epsilon]}(x_\epsilon) \to \text{Aut}_k(x_0)$
is surjective.
\end{enumerate}
\end{enumerate}
\end{lemma}

\begin{proof}
Assume $\mathcal{F}$ has (S1).
Suppose we have ring maps $f_i : A_i \to A$ in $\mathcal{C}_\Lambda$
with $f_2$ surjective. Let $x_i \in \mathcal{F}(A_i)$ such that
the pushforwards $f_{1, *}(x_1)$ and $f_{2, *}(x_2)$ are isomorphic.
Then we can denote $x$ an object of $\mathcal{F}$ over $A$ isomorphic
to both of these and we obtain a diagram as in (S1). Hence we find
an object $y$ of $\mathcal{F}$ over $A_1 \times_A A_2$ whose pushforward
to $A_1$, resp.\ $A_2$ is isomorphic to $x_1$, resp.\ $x_2$. In this way
we see that (S1) holds for $\overline{\mathcal{F}}$.

\medskip\noindent
Assume $\mathcal{F}$ has (S2).
The first part of (S2) for $\overline{\mathcal{F}}$ follows as in
the argument above. The second part of (S2) for
$\overline{\mathcal{F}}$ signifies that the map
$$
\overline{\mathcal{F}}(A \times_k k[\epsilon]) \to
\overline{\mathcal{F}}(A)
\times_{\overline{\mathcal{F}}(k)} \overline{\mathcal{F}}(k[\epsilon])
$$
is injective for any ring $A$ in $\mathcal{C}_\Lambda$. Suppose that
$y, y' \in \mathcal{F}(A \times_k k[\epsilon])$. Using the axioms
of cofibred categories we can choose commutative diagrams
$$
\vcenter{
\xymatrix{
y \ar[r]_c \ar[d]_a & x_\epsilon \ar[d]^e \\
x \ar[r]^d          & x_0
}
}
\quad\text{and}\quad
\vcenter{
\xymatrix{
y' \ar[r]_{c'} \ar[d]_{a'} & x'_\epsilon \ar[d]^{e'} \\
x' \ar[r]^{d'}                 & x'_0
}
}
\quad\text{lying over}\quad
\vcenter{
\xymatrix{
A \times_k k[\epsilon] \ar[d] \ar[r] & k[\epsilon] \ar[d] \\
A \ar[r] & k
}
}
$$
Assume that there exist isomorphisms
$\alpha : x \to x'$ in $\mathcal{F}(A)$ and
$\beta : x_\epsilon \to x'_\epsilon$ in $\mathcal{F}(k[\epsilon])$.
This also means there exists an isomorphism $\gamma : x_0 \to x'_0$
compatible with $\alpha$. To prove (S2) for $\overline{\mathcal{F}}$
we have to show that there exists an isomorphism $y \to y'$ in
$\mathcal{F}(A \times_k k[\epsilon])$.
By (S2) for $\mathcal{F}$ such a morphism will exist if we can
choose the isomorphisms $\alpha$ and $\beta$ and $\gamma$ such that
$$
\xymatrix{
x \ar[d]^\alpha \ar[r] & x_0 \ar[d]^\gamma &
x_\epsilon \ar[d]^\beta \ar[l]^e \\
x' \ar[r] & x'_0 & x'_\epsilon \ar[l]_{e'}
}
$$
is commutative (because then we can replace $x$ by $x'$ and $x_\epsilon$
by $x'_\epsilon$ in the previous displayed diagram). The left hand square
commutes by our choice of $\gamma$. We can factor $e' \circ \beta$ as
$\gamma' \circ e$ for some second map
$\gamma' : x_0 \to x'_0$. Now the question is whether we can arrange it so
that $\gamma = \gamma'$? This is clear if $\mathcal{F}(k)$ is a set, or a
setoid. Moreover, if
$\text{Aut}_{k[\epsilon]}(x_\epsilon) \to \text{Aut}_k(x_0)$
is surjective, then we can adjust the choice of $\beta$ by precomposing
with an automorphism of $x_\epsilon$ whose image is
$\gamma^{-1} \circ \gamma'$ to make things work.
\end{proof}

\end{document}
